\documentclass[10pt,a4paper]{amsart}
\usepackage[margin=19mm]{geometry}
\usepackage{amsmath,amssymb,mathtools}
\usepackage[scr=rsfs,cal=euler]{mathalfa}
\usepackage{xfrac}
\usepackage[unicode,hidelinks,bookmarks=false]{hyperref}
\newcommand{\ie}{i.e.\ }
\newcommand{\cf}{cf.\ }
\newcommand{\resp}{respectively\ }
\newcommand{\Subsection}{\subsection}
\providecommand{\nobreakdash}{\nobreak\hbox{-}}
\setlength{\emergencystretch}{2em}
\multlinegap0pt
\usepackage[shortlabels]{enumitem}
\usepackage{amssymb}
\newtheorem{theorem}{Theorem}
\newtheorem{lemma}{Lemma}
\newcommand{\K}{\mathrm{K}}      
\newcommand{\Osh}{{\mathcal O}}                        %  Structure sheaf
\newcommand{\PP}{\mathbb{P}} % projective space
\providecommand{\binom}[2]{{#1\choose#2}}
\title{Seshadri slope $\K$-semistability for the blow-up of projective space along a linear subspace}
\author{Nathan Grieve}
\date{Source: arXiv:2609.12319v1 (CC BY 4.0)}
\begin{document}
\maketitle

\noindent\emph{Source and licence.} Seshadri slope $\K$-semistability for the blow-up of projective space along a linear subspace. Version arXiv:2609.12319v1 (CC BY 4.0), distributed by arXiv under the Creative Commons Attribution 4.0 International licence, http://creativecommons.org/licenses/by/4.0/, as recorded in the arXiv licence metadata for this exact version and shown on the abstract page https://arxiv.org/abs/2609.12319v1. Full licence notice: \texttt{licenses/arxiv-2609.12319-CC-BY-4.0.txt}.\par\smallskip

\noindent\textit{Editorial scope.} Counted unit: the article's lemma slope:invariants:a:0 (source0.tex lines 486--499). The article's complete original proof of it is delivered with it (source0.tex lines 500--538, one \texttt{proof} environment of 39 lines). No mathematics is written, repaired or re-derived: the statement and the proof are the article's own lines, in the article's own notation, and no retained line is substituted or re-flowed.  Retained uncounted as the source-local context the proof needs: lines 461--466.  Nothing was omitted from any retained span, so the package carries no omission record.  No retained line contains a citation, so the wrapper carries no bibliography and no bibliography entry.  No retained line contains a figure environment, an image-inclusion command or drawing code, so no image is delivered and the package declares no asset dependency.  Editorial formatting only, in addition: an AMS \texttt{amsart} preamble that restates the article's own declarations and macros used by the retained lines, namely the environments \texttt{theorem}, \texttt{lemma} on separate counters, together with the standard packages those lines need.  The retained environments print in this document as Theorem 1, Lemma 1 in order of appearance, whereas the article prints the same items with the numbering of its own full document.  The complete native source of the article as distributed in the arXiv e-print for this exact version is bundled unchanged as \texttt{sources/210102/source0.tex}.
\begin{theorem}\label{instability:main:thm}
Let $\mathcal{E} = \Osh_{\PP^s}^{\oplus r} \oplus \Osh_{\PP^s}(1)$ and let $L_0 = \Osh_{\PP(\mathcal{E})}(1)$.  Then $L_0$ is Seshadri slope $\K$-semistable along the exceptional divisor $E$.  In more precise terms
\begin{equation}\label{slope:stab:key:eqn}
\mu_1(L_0;E) - \mu(L_0) = 0 \text{.}
\end{equation}
\end{theorem}

\begin{lemma}\label{slope:invariants:a:0}
In the setting of Theorem \ref{instability:main:thm} it holds true that
\begin{enumerate}
\item[(i)]{
$\int_0^1 \alpha_0(L_0;E,t) \mathrm{d}t = \frac{n-r+1}{(n+1)!}$
}
\item[(ii)]{
$ \int_0^1 \alpha_1(L_0;E,t) \mathrm{d}t = \frac{(n-r+1)(r+1)+(n-r)^2}{2n!}$
}
\item[(iii)]{
$\frac{1}{2} \int_0^1 \alpha_0'(L_0;E,t) \mathrm{d}t =  \frac{- 1}{2n!}$.
}
\end{enumerate}
\end{lemma}

\begin{proof}
For (i) note that
$$
\alpha_0(L_0;E,t) = \frac{1}{n!} \sum_{i=r}^n \binom{n}{i} (1-t)^i t^{n-i} \text{.}
$$
Thus
\begin{align*}
\int_0^1 \alpha_0(L_0;E,t)\mathrm{d}t &= \frac{1}{n!} \sum_{i=r}^n \frac{n!}{(n-i)!i!} \frac{i! (n-i)!}{(n+1)!} \\
& = \frac{1}{n!} \sum_{i=r}^n \frac{1}{n+1} 
= \frac{n-r+1}{(n+1)!} \text{.}
\end{align*}

For (ii)
\begin{align*}
\alpha_1(L_0;E) & = - \frac{1}{2(n-1)!} \left( -(r+1) \sum_{i=r-1}^{n-1} \binom{n-1}{i}(1-t)^i t^{n-1-i} \right. \\ 
& \left. - (n-r) \sum_{i=r}^{n-1}\binom{n-1}{i}(1-t)^it^{n-1-i}\right) \text{.}
\end{align*}
Thus
\begin{align*}
\int_0^1 \alpha_1(L_0;E,t)\mathrm{d}t &= - \frac{1}{2(n-1)!} \left( -(r+1) \sum_{i=r-1}^{n-1} \frac{(n-1)!}{i!(n-1-i)!} \frac{i!(n-1-i)!}{n!} \right. \\
& 
\left.  - (n-r) \sum_{i=r}^{n-1} \frac{ (n-1)!}{i!(n-1-i)!} \frac{i!(n-1-i)!}{n!} \right) \\
&= \frac{1}{2(n-1)!} \left( (r+1) \sum_{i=r-1}^{n-1} \frac{1}{n} + (n-r) \sum_{i=r}^{n-1} \frac{1}{n} \right) \\
& = \frac{1}{2n(n-1)!} \left( (r+1)(n-r+1) + (n-r)(n-r) \right) \\
& = \frac{ (n-r+1)(r+1) + (n-r)^2}{2n (n-1)! } \text{.}
\end{align*}

Finally, for (iii) 
\begin{align*}
\alpha_0'(L_0;E,t) &= \frac{-1}{(n-1)!} \left( \sum_{i=r-1}^{n-1} \binom{n-1}{i} (1 - t)^it^{n-1-i} \right. \\
& \left. - \sum_{i=r}^{n-1}\binom{n-1}{i} (1 - t)^i t^{n-1-i} \right) \text{.}
\end{align*}
Thus
\begin{align*}
\int_0^1 \frac{\alpha_0'(L_0;E,t)}{2} \mathrm{d}t &= \frac{-1}{2(n-1)!} \left( \sum_{i = r-1}^{n-1} \binom{n-1}{i} \frac{i!(n-1-i)!}{n!} \right. \\
& \left. - \sum_{i=r}^{n-1} \binom{n-1}{i} \frac{i!(n-1-i)!}{n!} \right) \\
&= \frac{-1}{2(n-1)!} \left( \frac{n-r+1}{n} - \frac{n-r}{n} \right) = \frac{-1}{2n(n-1)!} \text{.}
\end{align*}
\end{proof}

\end{document}
